\documentclass[../../main.tex]{subfiles}
\begin{document}

\tikzstyle{Default}=[circle, draw, thick]
\tikzstyle{Highlight}=[circle, draw, thick, red]

\begin{problem}
    给定一个树塔, 如下图所示.
    在此树塔中, 从顶部出发, 可以选择向左走还是向右走, 一直走到最低层. 
    请找出一条路径, 使得路径上的数值和最大, 并给出该数值和.
\end{problem}
\begin{center}
\begin{tikzpicture}
\node[Default] (A) at ( 0 , 0 ) {19};
\node[Default] (B) at ( 2 , 0 ) { 7};
\node[Default] (C) at ( 4 , 0 ) {10};
\node[Default] (D) at ( 6 , 0 ) { 4};
\node[Default] (E) at ( 8 , 0 ) {16};

\node[Default] (F) at ( 1, 1) { 2};
\node[Default] (G) at ( 3, 1) {18};
\node[Default] (H) at ( 5, 1) { 9};
\node[Default] (I) at ( 7, 1) { 5};

\node[Default] (J) at ( 2, 2) {10};
\node[Default] (K) at ( 4, 2) { 6};
\node[Default] (L) at ( 6, 2) { 8};

\node[Default] (M) at ( 3, 3) {12};
\node[Default] (N) at ( 5, 3) {15};

\node[Default] (O) at ( 4, 4) { 9};
\draw(F) -- (A);
\draw(F) -- (B);
\draw(G) -- (B);
\draw(G) -- (C);
\draw(H) -- (C);
\draw(H) -- (D);
\draw(I) -- (D);
\draw(I) -- (E);


\draw(J) -- (F);
\draw(J) -- (G);
\draw(K) -- (G);
\draw(K) -- (H);
\draw(L) -- (H);
\draw(L) -- (I);

\draw(M) -- (J);
\draw(M) -- (K);
\draw(N) -- (K);
\draw(N) -- (L);

\draw(O) -- (M);
\draw(O) -- (N);

\end{tikzpicture}
\end{center}

\begin{solution}
定义状态$dp[i][j]$表示从顶部到第$i$层, 选择第$j$个节点时的最大路径和.
注意到从顶点出发时到底是向左走还是向右走, 取决于向左走能取得最大值还是向右走能取得最大值.
因此, 可以得到状态转移方程如下:
\begin{equation*}
    dp[1][1] = max(dp[2][1], dp[2][2]) + 19
\end{equation*}
而对于\(dp[2][1]\), 也是同样的道理, 有:
\begin{equation*}
    dp[2][1] = max(dp[3][1], dp[3][2]) + 7
\end{equation*}
以此类推, 可以得到状态转移方程如下:
\begin{equation*}
    dp[i][j] = max(dp[i+1][j], dp[i+1][j+1]) + a[i][j]
\end{equation*}
在最后一层时, 可以直接得到最大路径和为$dp[5][j]$. 即:
\begin{equation*}
    dp[i][j] =
    \left\{
    \begin{aligned}
        & max(dp[i+1][j], dp[i+1][j+1]) + a[i][j] i \neq 5\\
        & a[i][j] , i=5
    \end{aligned}
    \right.
\end{equation*}
\(dp[i][j]\)对应的表格如下:
\begin{table}[ht!]
  \begin{center}
    \begin{tabular}{c|c|c|c|c|c} 
    \(dp[i][j]\) &  1 &  2 &  3 &  4 & 5 \\ \hline
    1            & 59 &  - &  - &  - & - \\ \hline
    2            & 50 & 49 &  - &  - & - \\ \hline
    3            & 38 & 34 & 29 &  - & - \\ \hline
    4            & 21 & 28 & 19 & 21 & - \\ \hline
    5            & 19 &  7 & 10 &  4 & 16\\ 
    \end{tabular}
  \end{center}
\end{table}
\newpage
容易得出最大路径和为$59$,其路径为:$9 \rightarrow 12 \rightarrow 10 \rightarrow 18 \rightarrow 10$.
即:
\begin{center}
\begin{tikzpicture}
\node[Default] (A) at ( 0 , 0 ) {19};
\node[Default] (B) at ( 2 , 0 ) { 7};
\node[Highlight] (C) at ( 4 , 0 ) {10};
\node[Default] (D) at ( 6 , 0 ) { 4};
\node[Default] (E) at ( 8 , 0 ) {16};

\node[Default] (F) at ( 1, 1) { 2};
\node[Highlight] (G) at ( 3, 1) {18};
\node[Default] (H) at ( 5, 1) { 9};
\node[Default] (I) at ( 7, 1) { 5};

\node[Highlight] (J) at ( 2, 2) {10};
\node[Default] (K) at ( 4, 2) { 6};
\node[Default] (L) at ( 6, 2) { 8};

\node[Highlight] (M) at ( 3, 3) {12};
\node[Default] (N) at ( 5, 3) {15};

\node[Highlight] (O) at ( 4, 4) { 9};
\draw(F) -- (A);
\draw(F) -- (B);
\draw(G) -- (B);
\draw(G) -- (C)[Highlight];
\draw(H) -- (C);
\draw(H) -- (D);
\draw(I) -- (D);
\draw(I) -- (E);


\draw(J) -- (F);
\draw(J) -- (G)[Highlight];
\draw(K) -- (G);
\draw(K) -- (H);
\draw(L) -- (H);
\draw(L) -- (I);

\draw(M) -- (J)[Highlight];
\draw(M) -- (K);
\draw(N) -- (K);
\draw(N) -- (L);

\draw(O) -- (M)[Highlight];
\draw(O) -- (N);

\end{tikzpicture}
\end{center}
\end{solution}
\end{document}